\documentclass[pdf]{beamer}
\mode<presentation>{}

\usepackage{multicol}
\usepackage{graphicx}
\usepackage{xcolor}
\usepackage{listings}
\usepackage{mips}
\usepackage{amsmath}
\usepackage{cancel}
\usepackage{tikz}
\usepackage{url}
\usetikzlibrary{shapes,backgrounds,arrows, fit, positioning}
\usepackage{algpseudocode}


\definecolor{CommentGreen}{rgb}{0,.6,0}
\lstset{
 language=[mips]Assembler,
 escapechar=@, % include LaTeX code between `@' characters
 keepspaces,   % needed to preserve spacing with lstinline
 basicstyle=\footnotesize\ttfamily\bfseries,
 commentstyle=\color{CommentGreen},
 stringstyle=\color{cyan},
 showstringspaces=false,
 keywordstyle=[1]\color{blue},    % instructions
 keywordstyle=[2]\color{magenta}, % directives
 keywordstyle=[3]\color{red},     % registers
 }
 
\lstset{language=[mips]Assembler}

%% preamble
\title{Tamarin: Concolic Disequivalence for MIPS}

\author{Abel Nieto}
\date{}

\newcommand{\crule}[3][black]{\textcolor{#1}{\rule{#2}{#3}}}

\begin{document}

\begin{frame}
\titlepage
\end{frame}

\begin{frame}{A Tale of Two Boxes}
~\\
\crule{3cm}{3cm} \hspace{2cm} \crule{3cm}{3cm}
\end{frame}

\begin{frame}{A Tale of Two Boxes}
\hspace{3.5cm}\fbox{\includegraphics[width=1cm]{carrot}}
\\
\crule{3cm}{3cm} \hspace{2cm} \crule{3cm}{3cm}
\pause
\\
nom nom
\pause
\hspace{3.5cm}
nom nom
\pause
\\
\hspace{3.5cm}
{\Huge $=$}
\end{frame}

\begin{frame}{A Tale of Two Boxes}
~\\
\fbox{\includegraphics[width=3cmcm]{camel}} \hspace{2cm} \fbox{\includegraphics[width=3cm]{beaver}}
\\
\hspace{5cm} {\Huge $\neq$}
\end{frame}

\Large

\begin{frame}{The Problem}
Given MIPS program $P_1$ and $P_2$, when are they equivalent?
\end{frame}

\begin{frame}{The Problem}
Attempt 1: two programs are equivalent if they give the same output (resp.) for all inputs.
\pause
\\~\\
\pause
What's an input?
\pause
Register $\$1$ and $\$2$.
\pause
\\
What's an output?
\pause
Register $\$3$. 
\pause
\\~\\
Don't care about (most) CPU interrupts/IO.
\end{frame}

\begin{frame}{The Problem}
Attempt 1: two programs are equivalent if they give the same output (resp.) for all inputs.
\pause
\\~\\
Problem: undecidable via Rice's theorem.
\end{frame}

\begin{frame}[fragile]{The Problem}
Attempt 2: two programs are $S$-equivalent if they cannot be told apart after $S$ steps.
\pause
\\~\\
$S$-equivalent (e.g. for $S = 10$), but not equivalent:
\begin{multicols}{2}
\begin{lstlisting}[escapeinside={(*}{*)}]
# P_1
add $3, $1, $2
\end{lstlisting}
\vfill\null
\columnbreak
\begin{lstlisting}
# P_2
  add $4, $0, 1  # counter
  add $5, $0, 42 # upper bound
loop:
  slt $6, $4, $5 
  beq $6, $0, end
  add $4, $4, 1
  beq $0, $0, loop
end:
  add $3, $1, $1  
\end{lstlisting}
\end{multicols}
\end{frame}

\begin{frame}{The Problem}
Attempt 2: two programs are $S$-equivalent if they \textbf{cannot be told apart} after $S$ steps.
\vspace{1cm}

\begin{tabular}{l | l | l}
$R_1$ & $R_2$ & $S$-equiv \\
\hline
\pause
$v$ & $v$ & yes \\
\pause
$v$ & $w \neq v$ & no \\
\pause
$v$ & error & no \\
\pause
error & error & yes \\
\pause
non-termination & ??? & yes \\
$\ldots$

\end{tabular}
\end{frame}

\begin{frame}{The Problem}
Attempt 2: two programs are $S$-equivalent if they \textbf{cannot be told apart} after $S$ steps.
\\~\\
Which inputs?
\pause 
\\~\\
Try some inputs by hand: low coverage, fast (unit tests)
\pause
\\~\\
Try all $2^{64}$values of $\$1$ and $\$2$: high coverage, slow (but decidable)
\pause
\\~\\
\textbf{Tamarin}: use concolic testing to check if the programs are disequivalent
\end{frame}

\end{document}